% page 204 of the facsimile (image p-203 of the scan)
% block 157.5 x 246.3 mm ; left margin 8.0 mm
\pgc{-12.700mm}{0.000mm}{
\l{\cel{4}196}
\l{}
\l{\sou{glikolo}. (\sou{kemio)}. Substanco sukroza, kristaligebla, quan onu}
\l{\cel{3}obtenas per duimigar (per saponigo) la acidi qui kompozas}
\l{\cel{3}bilo, o per traktar la gelatino per sulfacido. - DEFIRS.}
\l{}
\l{\sou{glikoproteido}.\cel{2}(\sou{kemio biologia}la)Asociuro molekulara ek pro-}
\l{\cel{3}teino e komplexo hidokarboza.}
\l{}
\l{\sou{glikoso}. (\sou{kemio)}. Elemento sukroza di la vito, di amilo, e c.}
\l{\cel{3}- DEFIRS.}
\l{}
\l{\sou{glikosido}. Nomo generala di la kompozaji naturala quin onu}
\l{\cel{3}renkontras che la vejetanti, e qui per hidrolizo, influe da}
\l{\cel{3}fermenti, produktas glikosi e substanci diversa-natura}
\l{\cel{3}(acidi, fenoli, alkoholi, e c.)}
\l{}
\l{\sou{glinar}.\cel{2}(\sou{trans.}) I. Rekoliar en agro (la spiki qui restas}
\l{\cel{3}pos la mesono). - II. (\sou{metaf}.)Rekoliar to quo lasesis da}
\l{\cel{3}altru. - EF.}
\l{}
\l{\sou{gliomo}. (\sou{patol.})\cel{2}Tumoro di la tisuo nervala.}
\l{}
\l{\sou{gliptar}. (\sou{trans}.) Skultar medaliie lapidi, kave (intalii) o}
\l{\cel{3}reliefe (kamei). - DEFIS.}
\l{}
\l{\sou{gliro}. (\sou{zool}.) Mikra mamifero rodera, qua torporas dum vintro.}
\l{\cel{3}- L.\cel{1}\sou{myoxus glis}.}
\l{}
\l{\sou{glitar}. (\sou{netrans}.)Movesar, movar, kontinue sur la surfaco di}
\l{\cel{3}korpo glata, konseque de pulso. - DEF.}
\l{}
\l{\sou{globo}. Korpo sfera o sferoida.\cel{2}- DEFIRS.}
\l{}
\l{\sou{globoido}.\cel{2}(\sou{bot}.) Korpo globatra od ovatra, qua ulakaze esas}
\l{\cel{3}inkluzita en grano di aleurono, e qua judikeble konsistas}
\l{\cel{3}ye glicerofosfato di kalko e di magnezio.}
\l{}
\l{\sou{globulo}. (\sou{anat.}) Korpuskulo plu o min ronda, suspendita en}
\l{\cel{3}sango. - EFIS.}
\l{}
\l{\sou{globulino}. (\sou{kemio biologiala)}. Materio albuminoida, qua}
\l{\cel{3}diferas de la albumini pro esar nedissolvebla en aquo dis-}
\l{\cel{3}tilita, ma dissolvebla en medii alkali-saloza.}
\l{}
\l{\sou{glorio}. Brilo de famo, de grandeso. - DEFIRS.}
\l{}
\l{\sou{gloso}. Expliko di la vorti di autoro literaturala, di qui la}
\l{\cel{3}senco divenis obskura pro lia obsoleteso. - DEFIS.}
\l{}
\l{\sou{glosar}io. I. Vortolibro kontenanta la vorti di qui la senco}
\l{\cel{3}divenis obskura pro lia obsoleteso, e pri qui expliko esas}
\l{\cel{3}nekareebla. - II. Nomenklaturo de la vorti ye qui konsistas}
\l{\cel{3}linguo. - DEF.}
\l{}
\l{\sou{gloto.(anat.}) Fenduro an la parto supra di laringo, qua uzesas}
\l{\cel{3}por la emiso di la voco. - EFIRS.}
}
